% ============================================================================
\section{Phương pháp}
\label{sec:method}
% ============================================================================

\subsection{Tổng quan}
\label{sec:overview}

Hình~\ref{fig:pipeline} là toàn bộ quy trình. Bốn hệ thống được chấm bằng \emph{cùng} một
harness đánh giá: hai baseline không học (luôn từ chối; TF-IDF) và hai Transformer được
fine-tune (PhoBERT, XLM-R). Điểm khác biệt duy nhất giữa hai nhánh Transformer là bước biến
văn bản thành token --- đúng biến mà câu hỏi nghiên cứu RQ1 muốn đo. Mọi hệ thống ghi ra dự
đoán \emph{từng câu}, để các phép đo chẩn đoán ở Mục~\ref{sec:alignment}--\ref{sec:paired}
ghép được với từng câu hỏi thay vì chỉ với con số tổng.

\begin{figure}[htbp]
\centering
\resizebox{\textwidth}{!}{%
\begin{tikzpicture}[font=\small,
  box/.style={draw=black!45,rounded corners=3pt,fill=white,align=center,minimum height=0.95cm,inner sep=4pt},
  pbox/.style={box,fill=uitblue!8,draw=uitblue},
  xbox/.style={box,fill=okgreen!8,draw=okgreen},
  dbox/.style={box,fill=orange!8,draw=warnorange},
  arr/.style={-{Latex[length=2mm]},black!60}]
  \node[box] (data) {UIT-ViQuAD 2.0\\\footnotesize train / validation};
  \node[pbox,right=0.8cm of data,yshift=0.85cm] (seg) {Tách từ \code{pyvi}\\\footnotesize ``Hà\_Nội''};
  \node[pbox,right=0.5cm of seg] (bpe) {BPE PhoBERT\\\footnotesize offset = khoảng của từ};
  \node[xbox,right=0.8cm of data,yshift=-0.85cm,minimum width=4.9cm] (spm) {SentencePiece XLM-R\\\footnotesize offset từ fast tokenizer};
  \node[box,right=0.6cm of bpe,yshift=-0.85cm] (win) {Cửa sổ trượt\\\footnotesize offset \emph{tuyệt đối}};
  \node[box,right=0.5cm of win] (enc) {Encoder + lớp QA\\\footnotesize $(S_i, E_i)$};
  \node[box,right=0.5cm of enc] (dec) {Chọn span / từ chối\\\footnotesize cắt từ context gốc};
  \node[dbox,below=0.9cm of dec] (pred) {Dự đoán\\từng câu};
  \node[dbox,left=0.5cm of pred] (met) {EM / F1\\\footnotesize tách có/không đáp án};
  \node[dbox,left=0.5cm of met,align=left] (diag) {\footnotesize$\bullet$ biên từ + trần EM\\\footnotesize$\bullet$ phân loại lỗi\\\footnotesize$\bullet$ precision/recall từ chối\\\footnotesize$\bullet$ McNemar + bootstrap};
  \draw[arr] (data.east) -- ++(0.3,0) |- (seg.west);
  \draw[arr] (data.east) -- ++(0.3,0) |- (spm.west);
  \draw[arr] (seg) -- (bpe);
  \draw[arr] (bpe.east) -| ([xshift=-0.3cm]win.north);
  \draw[arr] (spm.east) -| ([xshift=-0.3cm]win.south);
  \draw[arr] (win) -- (enc); \draw[arr] (enc) -- (dec); \draw[arr] (dec) -- (pred);
  \draw[arr] (pred) -- (met); \draw[arr] (met) -- (diag);
\end{tikzpicture}}
\caption{Quy trình. Nhánh xanh dương (PhoBERT) và xanh lá (XLM-R) chỉ khác nhau ở bước tạo token và offset; mọi bước sau dùng chung mã}
\label{fig:pipeline}
\end{figure}

\paragraph{Tái sử dụng mã nguồn.} Tầng dữ liệu, metric, cửa sổ trượt, vòng lặp huấn luyện và
harness đánh giá (\code{src/mrc/}) được kế thừa từ đồ án môn CS116 của cùng nhóm tác giả, vốn
đã có bộ kiểm thử. Đồ án CS221 bổ sung phần phục vụ câu hỏi nghiên cứu riêng: lớp tokenizer
cho PhoBERT, phép đo biên từ, kiểm định stress-test, các phép đo chẩn đoán và giao thức chọn
epoch trên tập dev tách từ train. Hai đồ án trả lời hai câu hỏi khác nhau trên hai tập mô hình
khác nhau; không con số nào của đồ án CS116 được dùng lại trong báo cáo này. Khung mã ban đầu
của chính đồ án CS221 (thư mục \code{legacy/}) \emph{không} được dùng để sinh kết quả vì các
lỗi được nêu ở Phụ lục~\ref{app:legacy}.

\subsection{Giao thức dữ liệu}
\label{sec:protocol}

\begin{itemize}
  \item \textbf{Test chính thức của UIT-ViQuAD~2.0 là tập ẩn nhãn.} Toàn bộ 7.301 câu có đáp
        án rỗng nhưng \code{is\_impossible = False}; không chấm được. Harness từ chối chấm các
        câu như vậy bằng một kiểm tra chặn (\code{assert\_gradeable}).
  \item \textbf{Validation chính thức (\NVal{} câu) đóng vai tập test của đồ án}, được chấm
        \emph{đúng một lần} cho mỗi hệ thống, trên toàn bộ tập, không lấy mẫu con.
  \item \textbf{Chọn epoch trên tập dev tách từ train}: \NDevCtx{} đoạn văn (\NDev{} câu, 5\%
        số đoạn văn) được tách khỏi train \emph{theo đoạn văn}, để không câu hỏi nào của cùng
        một đoạn văn nằm ở cả hai phía. Mô hình được huấn luyện trên phần còn lại (\NFit{} câu).
        Như vậy validation không ảnh hưởng tới bất kỳ quyết định nào trước khi chấm.
  \item \textbf{Chống rò rỉ}: trước khi huấn luyện, chương trình khẳng định (\code{assert})
        rằng ba tập train/dev/validation không có đoạn văn chung; vi phạm làm dừng chương trình.
  \item \textbf{Bộ stress-test chỉ dùng để đánh giá}, không bao giờ trộn vào dữ liệu huấn luyện
        (khung mã cũ làm điều này --- Phụ lục~\ref{app:legacy}).
\end{itemize}

\subsection{Hai baseline}
\label{sec:baselines}

\paragraph{Luôn từ chối.} Trả về chuỗi rỗng cho mọi câu. Theo quy ước chấm SQuAD~2.0, hệ thống
này đạt điểm bằng đúng tỉ lệ câu không có đáp án của tập chấm. Nó không phải một mô hình mà là
một \emph{thước đo}: mọi con số khác phải được đọc so với mức sàn này, và nó làm lộ ra những
nhóm câu hỏi mà điểm cao không chứng minh năng lực gì (Mục~\ref{sec:res-rq3}).

\paragraph{TF-IDF.} Truy hồi câu bằng TF-IDF (Mục~\ref{sec:tfidf}): tách đoạn văn thành câu,
vector hoá câu hỏi và các câu bằng \code{TfidfVectorizer} (n-gram 1--2), trả về câu có cosine
lớn nhất. Baseline này không bao giờ từ chối và luôn trả về \emph{cả một câu}.

\subsection{Đưa PhoBERT vào bài toán trích xuất}
\label{sec:segtok}

PhoBERT cần đầu vào đã tách từ nhưng không cung cấp offset. Lớp \code{SegmentedTokenizer} tự
dựng offset theo Thuật toán~\ref{alg:segtok}: tách từ bằng \code{pyvi}, căn từng từ về khoảng
ký tự của nó trong văn bản gốc, rồi mã hoá BPE \emph{từng từ}; mọi mảnh BPE của một từ mang
chung khoảng ký tự của cả từ. Lớp này cài đúng phần giao diện mà mã cửa sổ trượt và mã suy luận
sử dụng, nên PhoBERT và XLM-R đi qua \emph{cùng} một đường mã từ đó trở đi.

\begin{algorithm}[htbp]
\caption{Dựng offset cho PhoBERT từ bước tách từ}
\label{alg:segtok}
\KwIn{văn bản gốc $T$; bộ tách từ \textsc{Seg}; tokenizer BPE}
\KwOut{danh sách \code{ids}, danh sách \code{offsets} (khoảng ký tự trong $T$)}
$p \leftarrow 0$\;
\ForEach{từ $w$ trong \textsc{Seg}$(T)$}{
  $s \leftarrow \bot$\;
  \ForEach{âm tiết $a$ trong $w$ tách theo ``\_''}{
    $p \leftarrow$ vị trí đầu tiên $\ge p$ trong $T$ bắt đầu bằng $a$, bỏ qua khoảng trắng\;
    \If{không tìm thấy}{\Continue\;}
    \If{$s = \bot$}{$s \leftarrow p$\;}
    $p \leftarrow p + |a|$\;
  }
  \ForEach{mảnh BPE $b$ của $w$}{thêm $\mathrm{id}(b)$ vào \code{ids}; thêm $(s, p)$ vào \code{offsets}\;}
}
\Return \code{ids}, \code{offsets}\;
\end{algorithm}

Hai hệ quả của thiết kế này được dùng trực tiếp trong phân tích:

\begin{enumerate}
  \item \textbf{Đáp án luôn được cắt từ văn bản gốc}, nên không bao giờ chứa dấu gạch dưới
        của bước tách từ. Cách ngây thơ --- giải mã token bằng \code{decode()} --- trả về
        ``Hà\_Nội''; vì bước chuẩn hoá của metric xoá dấu câu, chuỗi này thành ``hànội'' và
        một đáp án \emph{đúng} bị chấm EM~0, F1~0. Lỗi này tồn tại trong khung mã ban đầu và
        được kiểm thử hồi quy.
  \item \textbf{PhoBERT chỉ trả lời được bằng dãy từ nguyên vẹn} --- đúng tính chất mà phép
        đo ở Mục~\ref{sec:alignment} khai thác.
\end{enumerate}

Tính đúng đắn được kiểm thử bằng các bất biến: mỗi offset cắt ra đúng từ gốc (thay ``\_''
bằng khoảng trắng); nhãn huấn luyện được giải mã ra lại đúng đáp án vàng khi đáp án khớp biên
từ; và bị \emph{mở rộng} đúng tới từ bao quanh khi đáp án cắt ngang từ. Trên validation, các từ
đã căn phủ \SegCtxCoverage\% ký tự không phải khoảng trắng của mọi đoạn văn --- không ký tự nào
bị bỏ sót.

\paragraph{Bộ tách từ.} PhoBERT được tiền huấn luyện với RDRSegmenter (VnCoreNLP), cần Java.
Đồ án dùng \code{pyvi}, thuần Python, để pipeline chạy được trên máy không có Java. Hai bộ tách
từ không trùng nhau hoàn toàn; đây là một nguồn lệch phân phối được nêu trong giới hạn
(Mục~\ref{sec:limitations}) và là một biến đáng thử trong hướng phát triển.

\subsection{Fine-tuning}
\label{sec:finetune}

Hai mô hình được huấn luyện bằng \emph{cùng} vòng lặp, cùng lịch học và cùng seed
(Bảng~\ref{tab:hparams}). Cần tách rõ hai loại khác biệt. PhoBERT chỉ có 256 vị trí hữu dụng nên
\code{max\_length} = 256 là \emph{bắt buộc}. Ngược lại, việc cho XLM-R dùng 384 token và chồng lấp
128 là \emph{lựa chọn} của nhóm (giá trị mặc định cho SQuAD), không phải ràng buộc kiến trúc --- XLM-R
hoàn toàn chạy được ở 256/96. Vì mô hình trả lời khi \emph{có} một cửa sổ vượt ngưỡng, số cửa sổ
khác nhau cũng ảnh hưởng nhẹ tới hành vi từ chối; Mục~\ref{sec:seeds} đo lại XLM-R ở đúng 256/96. Độ dài đáp án tối đa được đặt 64 token cho cả hai, lớn hơn phân vị
95 của độ dài đáp án vàng theo cả hai tokenizer (đo trên 400 câu train: 30 token PhoBERT, 38
token XLM-R), để tham số này không trở thành một yếu tố gây nhiễu.

\begin{table}[htbp]
\centering
\caption{Siêu tham số huấn luyện}
\label{tab:hparams}
\renewcommand{\arraystretch}{1.13}
\small
\begin{tabular}{@{}l c c@{}}
\toprule
\textbf{Tham số} & \textbf{PhoBERT-base-v2} & \textbf{XLM-R-base} \\
\midrule
Checkpoint & \code{vinai/phobert-base-v2} & \code{FacebookAI/xlm-roberta-base} \\
Tiền xử lý & tách từ \code{pyvi} & không \\
\code{max\_length} / chồng lấp cửa sổ & 256 / 96 & 384 / 128 \\
Độ dài đáp án tối đa (token) & 64 & 64 \\
Số epoch tối đa & 3 & 3 \\
Batch $\times$ tích luỹ gradient & $12 \times 2 = 24$ & $12 \times 2 = 24$ \\
Tốc độ học, warmup, weight decay & $3\cdot10^{-5}$; 10\%; 0,01 & $3\cdot10^{-5}$; 10\%; 0,01 \\
Tối ưu, cắt gradient & AdamW; 1,0 & AdamW; 1,0 \\
Chọn epoch & F1 cao nhất trên dev (500 câu) & F1 cao nhất trên dev (500 câu) \\
Ngưỡng không trả lời $\tau$ & 0; và chọn trên dev (Mục~\ref{sec:threshold}) & 0; và chọn trên dev \\
Seed & 42 & 42 \\
\bottomrule
\end{tabular}
\par\vspace{2pt}
{\footnotesize Nguồn: trường \code{config} trong \code{results/training\_curve\_\{phobert,xlmr\}.json}.}
\end{table}

Không có tìm kiếm siêu tham số: tốc độ học $3\cdot10^{-5}$ và lịch học là giá trị phổ biến cho
mô hình cỡ \emph{base} trên SQuAD. Đây là lựa chọn có chủ đích --- với một lần chạy mỗi mô hình,
tìm kiếm siêu tham số trên dev rồi báo cáo trên cùng dev sẽ thổi phồng kết quả --- nhưng cũng có
nghĩa là cả hai mô hình có thể chưa đạt điểm tốt nhất của chúng.

\subsection{Phép đo tương thích biên từ}
\label{sec:alignment}

Đây là công cụ chẩn đoán trung tâm cho RQ1, và nó \emph{không cần chạy mô hình}. Gọi
$\mathcal{W} = \{(a_k, b_k)\}$ là tập khoảng ký tự của các từ mà bộ tách từ sinh ra trên đoạn
văn, và $[s, e)$ là khoảng ký tự của một đáp án vàng. Đáp án \textbf{khớp biên từ}
(\emph{aligned}) nếu không đầu nào của nó nằm \emph{bên trong} một từ:
\begin{equation}
  \mathrm{aligned}(s, e) \iff \nexists\, (a, b) \in \mathcal{W}:\; a < s < b \;\lor\; a < e < b .
\end{equation}
Nếu không khớp, span nguyên-từ tốt nhất mà một mô hình word-level trả về được là span nhỏ nhất
bao đáp án:
\begin{equation}
  \hat{s} = \min\{a : (a,b)\in\mathcal{W},\, a < e,\, b > s\}, \qquad
  \hat{e} = \max\{b : (a,b)\in\mathcal{W},\, a < e,\, b > s\},
\end{equation}
và \textbf{trần EM} của câu đó là $\max_{g \in G} \mathrm{EM}(C[\hat{s}{:}\hat{e}], g)$. Với câu
có nhiều đáp án vàng, câu được tính là khớp biên nếu \emph{ít nhất một} đáp án khớp. Trung bình
trần EM trên toàn tập là điểm EM cao nhất mà một mô hình word-level \emph{hoàn hảo} có thể đạt,
chỉ tính riêng giới hạn do đơn vị đầu ra.

Phép đo này tách được hai kênh mà giả thuyết RQ1 gộp làm một: (i) kênh \emph{đơn vị đầu ra} ---
PhoBERT không thể trả lời nửa từ, đo được chính xác và không cần mô hình; và (ii) kênh \emph{biểu
diễn} --- lỗi tách từ làm sai biểu diễn đầu vào, chỉ đo được gián tiếp qua hiệu năng. Nếu kênh (i)
chỉ giải thích được một phần nhỏ khoảng cách giữa hai mô hình, phần còn lại phải đến từ nơi khác.

\subsection{Phân loại lỗi tự động}
\label{sec:taxonomy}

Mỗi dự đoán sai được xếp vào đúng một trong sáu loại (Bảng~\ref{tab:taxonomy-def}), định nghĩa
bằng quan hệ giữa đa tập token của dự đoán và của đáp án vàng sau chuẩn hoá. Với nhiều đáp án
vàng, dùng đáp án có F1 cao nhất với dự đoán. Vì định nghĩa không cần người gán, phép phân loại
chạy được trên \emph{toàn bộ} tập chấm và tái lập được.

\begin{table}[htbp]
\centering
\caption{Sáu loại lỗi}
\label{tab:taxonomy-def}
\renewcommand{\arraystretch}{1.15}
\small
\begin{tabularx}{\textwidth}{@{}l L L@{}}
\toprule
\textbf{Loại} & \textbf{Điều kiện} & \textbf{Ví dụ (vàng $\to$ dự đoán)} \\
\midrule
Từ chối sai & Có đáp án, dự đoán rỗng & ``Hà Nội'' $\to$ $\varnothing$ \\
Trả lời sai & Không có đáp án, dự đoán khác rỗng & $\varnothing$ $\to$ ``năm 1945'' \\
Biên thừa & Dự đoán \emph{chứa} đáp án vàng & ``Hà Nội'' $\to$ ``thành phố Hà Nội'' \\
Biên thiếu & Dự đoán \emph{nằm trong} đáp án vàng & ``thành phố Hà Nội'' $\to$ ``Hà Nội'' \\
Biên lệch & Có token chung, lệch cả hai phía & ``Hà Nội là'' $\to$ ``Nội là thủ đô'' \\
Sai vị trí & Không có token chung & ``Hà Nội'' $\to$ ``Paris'' \\
\bottomrule
\end{tabularx}
\end{table}

Ba loại ``biên'' là lỗi \emph{đã tìm đúng vùng nhưng sai ranh giới} --- loại lỗi mà giả thuyết RQ1
dự đoán PhoBERT mắc nhiều hơn. ``Sai vị trí'' là lỗi hiểu; ``từ chối sai'' và ``trả lời sai'' là lỗi
của quyết định có trả lời hay không.

\subsection{Hành vi từ chối như một bộ phân loại}
\label{sec:abstention}

Trên SQuAD~2.0, điểm tổng trộn hai kỹ năng. Để tách kỹ năng thứ hai, việc ``trả về rỗng'' được
coi là một bộ phân loại nhị phân dự đoán ``câu này không có đáp án'', với precision (trong những
lần từ chối, bao nhiêu lần đúng) và recall (trong những câu không có đáp án, bao nhiêu câu được
từ chối). Hệ thống luôn từ chối có recall 100\% và precision bằng tỉ lệ câu không có đáp án ---
một mô hình ``biết từ chối'' thật phải có precision cao hơn hẳn mức đó.

\subsection{So sánh cặp PhoBERT và XLM-R}
\label{sec:paired}

Chênh lệch EM giữa hai mô hình được kiểm định bằng McNemar chính xác (công thức~\ref{eq:mcnemar})
trên các câu mà hai mô hình bất đồng, kèm khoảng tin cậy bootstrap 95\% (2.000 lần lấy mẫu, seed
42) --- trên toàn tập và trên từng nhóm: câu có đáp án, câu không có đáp án, câu lệch biên từ.
Kết luận ``A khác B'' chỉ được đưa ra khi khoảng tin cậy không chứa 0; ngược lại báo cáo ghi
``chưa đủ bằng chứng''.

\subsection{Kiểm định bộ stress-test}
\label{sec:audit-method}

Trước khi dùng bộ stress-test 250 câu (5 nhóm E1--E5 $\times$ 50) để chấm mô hình, mỗi câu được
kiểm tra tự động: (1) với câu có đáp án, đáp án vàng có nằm trong đoạn văn không; (2)
\code{answer\_start} có trỏ đúng vào đáp án không; (3) đáp án có xuất hiện nguyên văn trong câu
hỏi không (rò rỉ đáp án); (4) số câu hỏi khác nhau và số khung câu hỏi (câu hỏi sau khi thay
cụm trong nháy đơn bằng \code{X}). Một câu \textbf{hợp lệ} nếu nó không có đáp án, hoặc đáp án của
nó nằm trong đoạn văn --- điều kiện tối thiểu để một mô hình trích xuất \emph{có thể} trả lời
đúng. Kết quả ở Mục~\ref{sec:audit}.

\subsection{Truy vết số liệu và kiểm thử}
\label{sec:traceability}

Mọi hệ thống ghi \code{results/eval\_<hệ~thống>\_<tập>.json} (số tổng, tách nhóm, commit, thời
điểm, thiết bị) và \code{results/predictions\_<hệ~thống>\_<tập>.json} (dự đoán từng câu). Mọi
phép đo chẩn đoán đọc các tệp này. Báo cáo \emph{không gõ tay con số kết quả nào}: script
\code{make\_report\_numbers.py} sinh các macro \LaTeX{} và các bảng từ \code{results/}; giá trị
nào chưa có sẽ hiện chữ đỏ ``??'' trong PDF thay vì một con số đoán. Mã nguồn có 244 kiểm thử
tự động chạy không cần GPU, gồm các bất biến của lớp tokenizer PhoBERT, phép đo biên từ, phân
loại lỗi và kiểm định McNemar.
